\section{Discussion}

\subsection{Main finding}
Under a strictly participant-wise, pre-specified protocol, task EEG from the MODMA dot-probe recordings did not
separate MDD from healthy controls better than chance. The conclusion holds across all model types tested:
\begin{itemize}
  \item the pre-specified primary model and every feature family on its own;
  \item linear, kernel and tree-based classifiers;
  \item Riemannian models and three deep networks trained on the full 128-channel epochs.
\end{itemize}
The one nominally significant secondary result (L1-regularised logistic regression) did not survive correction for
the number of secondary models. The honest estimate of picking the best model, obtained by nested model selection,
fell to \NestedmodelselectionAllAuc, with a confidence interval including chance.

By contrast, three demographic variables (age, sex and education) reached AUC \DemographicsLrAuc. Any EEG model that
claims diagnostic value on this dataset should be benchmarked against this figure, not only against 0.5.

These findings agree with reviews describing EEG markers of depression as small and heterogeneous
effects~\cite{olbrich2013eeg,deaguiar2019review,simmatis2023technical}. They also agree with the meta-analytic null
result for frontal alpha asymmetry~\cite{vandervinne2017faa}. They do not agree with the near-perfect accuracies often
reported for EEG-based depression detection~\cite{ay2019automated}.

\subsection{Why do published accuracies differ so much?}
Our inflation control offers a direct explanation. On the same data, the tangent-space model evaluated with a
trial-wise split ranked every participant correctly (participant-level AUC \InflTsLrTrialSplitSubjectAuc;
trial-level \InflTsLrTrialSplitTrialAuc). Evaluated on held-out participants, it was at chance
(\InflTsLrSubjectSplitSubjectAuc).

When epochs of one person appear in both training and test sets, the classifier only needs to recognise the person,
and every person carries one label. Individual EEG signatures are highly distinctive, so this task is easy, and it says
nothing about diagnosing a new patient~\cite{saeb2017need,brookshire2024leakage}. Reported accuracies therefore cannot
be interpreted unless the split is stated to be across participants. In this light, values above 95\% on datasets of
a few dozen participants should be regarded with suspicion.

\subsection{Model selection and multiplicity}
With $n=\NIncluded$, sampling noise alone produces large AUCs: under label permutation, the primary pipeline exceeded
AUC \PrimaryLrLAllNullQ{} in 5\% of runs. Across 18 configurations, some will exceed that by chance. Reporting the best of them, here L1-LR at AUC \LrLAllAuc, would have yielded a publishable but
misleading result.

Three elements of the design guard against this:
\begin{itemize}
  \item a primary model and outcome fixed before any result existed;
  \item multiplicity correction across secondary models;
  \item a nested estimate of the selection procedure itself~\cite{cawley2010overfitting}.
\end{itemize}
We recommend all three as standard practice for small clinical EEG datasets.

\subsection{Confounds}
The groups differed in education and, partly, in recording period: some months contain only patients or only
controls. In principle, period effects (equipment state, staff, season) could produce group differences unrelated to
depression.

Two observations make this unlikely to be the explanation here. First, the EEG features did not predict recording
period above chance. Second, the analysis restricted to the months in which both groups were recorded was no more
successful than the full analysis.

The signal-quality measures did not differ between groups. Even so, one descriptive pattern deserves caution:
\begin{itemize}
  \item MDD participants showed numerically higher beta and low-gamma power, strongest at peripheral electrodes;
  \item they also showed higher entropy and Lempel--Ziv complexity, the family whose permutation importance was
        largest, though not significant.
\end{itemize}
Both patterns are also typical signatures of residual muscle activity. ICA with ICLabel removes the clearest muscle
components but not all EMG. Any future report of high-frequency or complexity differences in these data should
therefore rule out muscle tension before attributing them to depression.

The single connectivity feature that survived FDR correction (left central--temporal theta wPLI) was not robustly
used by any model. It should be considered a hypothesis for replication, not a finding.

\subsection{Limitations}
\begin{itemize}
  \item \textbf{Sample size.} With \NIncludedMDD{} MDD and \NIncludedHC{} HC participants after exclusion, the
        participant-level confidence intervals are wide: for the primary model,
        \PrimaryLrLAllAucLo--\PrimaryLrLAllAucHi. A modest true effect, around AUC 0.6, cannot be excluded. Our
        conclusion is that these data provide no evidence of a diagnostic signal, not that none exists.
  \item \textbf{Single site and single task.} All data come from one site, one EEG system and one paradigm. Clinical
        details that may modulate EEG are not modelled, notably medication status and illness duration.
  \item \textbf{Asymmetric exclusions.} The pre-specified quality criteria excluded more patients than controls. The
        analysis on all 53 participants, however, gave the same conclusion.
  \item \textbf{Deep learning budget.} For computational reasons, the deep networks were evaluated with 3 rather than
        10 repetitions and without permutation tests. Their confidence intervals all included chance, and their point
        estimates were below that of the primary model.
  \item \textbf{Feature space.} The features were fixed a priori from the literature. Other representations could
        perform better, but searching for them on the same 45 participants would reintroduce the selection bias this
        study was designed to avoid.
\end{itemize}

\subsection{Conclusion}
Evaluated rigorously on participants never seen by the model, 128-channel task EEG from the MODMA dot-probe dataset
did not distinguish major depression from healthy controls, and was outperformed by a demographic baseline. The large
gap between trial-wise and participant-wise estimates on the same data shows how evaluation choices can manufacture
apparent diagnostic accuracy. Credible progress towards EEG markers of depression will require participant-wise
evaluation, pre-specified analyses, demographic baselines and substantially larger, multi-site samples.
